% GENERATED FILE. Edit canonical lessons, then run npm run generate:books.
% canonical-source-hash: fnv1a64:c400233ffbbcf8ca
% canonical-lessons: MR-C508-vetan, MR-C508-majuri, MR-C508-mandhan, MR-C508-pagarvadh, MR-C508-kamavne

\chapter{Wages and Earnings}
\label{ch:mr-wages-and-earnings}
\hlchaptermodality{508}

\begin{chapteropening}
\textbf{By the end of this chapter:} \emph{I can talk about pay, wages and earning a living.}

Complete the last lesson of chapter 508: I can talk about pay, wages and earning a living.
\end{chapteropening}

\section[\texorpdfstring{\mr{वेतन} (vetan)}{वेतन (vetan)}]{\mr{वेतन} (vetan) — a wage; pay}

\label{lesson:MR-C508-vetan}



\begin{quote}
Before the new one: say the Marathi for a knack, then the Marathi for to join a post.
\end{quote}

\subsection*{You'll want to know: \mr{वेतन}}
\textbf{\mr{वेतन}} — \emph{vetan} — \textquotedblleft{}a wage; pay\textquotedblright{}.

\begin{quote}
\mr{या} \mr{महिन्याचं} \mr{वेतन} \mr{उशिरा} \mr{आलं}.
\end{quote}

\begin{quote}
\emph{yā mahinyācaṁ vetan uśirā ālaṁ.}
\end{quote}

\begin{quote}
This month's pay came late.
\end{quote}

Neuter.

\mr{वेतन} is the formal word beside \mr{पगार}.

\begin{cousinweb}
\textbf{\mr{वेतन}} (\emph{vetan}) is a tatsama, taken from Sanskrit \emph{vetana}, \textquotedblleft{}wages, hire, salary\textquotedblright{} (Monier-Williams).
\end{cousinweb}

\begin{grammarlens}[title={What you've built}]
One more word for describing what you lived and felt.
\end{grammarlens}

\subsection*{Your turn}
\begin{itemize}
\raggedright
  \item \emph{Say it:} \emph{vetan}
  \item \emph{Say it:} \emph{vetan}, once more
  \item \emph{Say it:} \emph{yā mahinyācaṁ vetan uśirā ālaṁ.} — This month's pay came late.
  \item \emph{Recall:} say the Marathi for a knack, then the Marathi for to join a post, then say \emph{vetan} again
\end{itemize}

\begin{tcolorbox}[breakable,colback=teal!4,colframe=teal!35!black,title={Before you move on}]
What does \mr{वेतन} mean? (\textquotedblleft{}A wage; pay\textquotedblright{}.) Say it once more.
\end{tcolorbox}
\section[\texorpdfstring{\mr{मजुरी} (majurī)}{मजुरी (majurī)}]{\mr{मजुरी} (majurī) — daily wages}

\label{lesson:MR-C508-majuri}



\begin{quote}
Before the new one: say the Marathi for to join a post, then the Marathi for a wage.
\end{quote}

\subsection*{You'll want to know: \mr{मजुरी}}
\textbf{\mr{मजुरी}} — \emph{majurī} — \textquotedblleft{}daily wages\textquotedblright{}.

\begin{quote}
\mr{कामगारांना} \mr{संध्याकाळी} \mr{त्यांची} \mr{मजुरी} \mr{मिळाली}.
\end{quote}

\begin{quote}
\emph{kāmgārānnā sandhyākāḷī tyāñcī majurī miḷālī.}
\end{quote}

\begin{quote}
The workers got their wages in the evening.
\end{quote}

Feminine.

\begin{cousinweb}
\textbf{\mr{मजुरी}} (\emph{majurī}) is a loan from Persian \emph{muzdūrī}, \textquotedblleft{}work for hire\textquotedblright{}, built on \emph{muzd}, \textquotedblleft{}wage\textquotedblright{}.
\end{cousinweb}

\begin{grammarlens}[title={What you've built}]
One more word for describing what you lived and felt.
\end{grammarlens}

\subsection*{Your turn}
\begin{itemize}
\raggedright
  \item \emph{Say it:} \emph{majurī}
  \item \emph{Say it:} \emph{majurī}, once more
  \item \emph{Say it:} \emph{kāmgārānnā sandhyākāḷī tyāñcī majurī miḷālī.} — The workers got their wages in the evening.
  \item \emph{Recall:} say the Marathi for to join a post, then the Marathi for a wage, then say \emph{majurī} again
\end{itemize}

\begin{tcolorbox}[breakable,colback=teal!4,colframe=teal!35!black,title={Before you move on}]
What does \mr{मजुरी} mean? (\textquotedblleft{}Daily wages\textquotedblright{}.) Say it once more.
\end{tcolorbox}
\section[\texorpdfstring{\mr{मानधन} (māndhan)}{मानधन (māndhan)}]{\mr{मानधन} (māndhan) — an honorarium; a fee}

\label{lesson:MR-C508-mandhan}



\begin{quote}
Before the new one: say the Marathi for a wage, then the Marathi for daily wages.
\end{quote}

\subsection*{You'll want to know: \mr{मानधन}}
\textbf{\mr{मानधन}} — \emph{māndhan} — \textquotedblleft{}an honorarium; a fee\textquotedblright{}.

\begin{quote}
\mr{भाषणासाठी} \mr{लेखकाला} \mr{मानधन} \mr{मिळालं}.
\end{quote}

\begin{quote}
\emph{bhāṣṇāsāṭhī lekhakālā māndhan miḷālaṁ.}
\end{quote}

\begin{quote}
The writer got a fee for the talk.
\end{quote}

Neuter.

\begin{cousinweb}
\textbf{\mr{मानधन}} (\emph{māndhan}) is built from two tatsamas, Sanskrit \emph{māna}, \textquotedblleft{}honour\textquotedblright{}, and \emph{dhana}, \textquotedblleft{}wealth\textquotedblright{} (Monier-Williams): money given as a mark of respect.
\end{cousinweb}

\begin{grammarlens}[title={What you've built}]
One more word for describing what you lived and felt.
\end{grammarlens}

\subsection*{Your turn}
\begin{itemize}
\raggedright
  \item \emph{Say it:} \emph{māndhan}
  \item \emph{Say it:} \emph{māndhan}, once more
  \item \emph{Say it:} \emph{bhāṣṇāsāṭhī lekhakālā māndhan miḷālaṁ.} — The writer got a fee for the talk.
  \item \emph{Recall:} say the Marathi for a wage, then the Marathi for daily wages, then say \emph{māndhan} again
\end{itemize}

\begin{tcolorbox}[breakable,colback=teal!4,colframe=teal!35!black,title={Before you move on}]
What does \mr{मानधन} mean? (\textquotedblleft{}An honorarium; a fee\textquotedblright{}.) Say it once more.
\end{tcolorbox}
\section[\texorpdfstring{\mr{पगारवाढ} (pagārvāḍh)}{पगारवाढ (pagārvāḍh)}]{\mr{पगारवाढ} (pagārvāḍh) — a pay rise}

\label{lesson:MR-C508-pagarvadh}



\begin{quote}
Before the new one: say the Marathi for daily wages, then the Marathi for an honorarium.
\end{quote}

\subsection*{You'll want to know: \mr{पगारवाढ}}
\textbf{\mr{पगारवाढ}} — \emph{pagārvāḍh} — \textquotedblleft{}a pay rise\textquotedblright{}.

\begin{quote}
\mr{या} \mr{वर्षी} \mr{सगळ्यांना} \mr{चांगली} \mr{पगारवाढ} \mr{मिळाली}.
\end{quote}

\begin{quote}
\emph{yā varṣī sagḷyānnā cāṅglī pagārvāḍh miḷālī.}
\end{quote}

\begin{quote}
Everyone got a good pay rise this year.
\end{quote}

Feminine.

\mr{पगार} is masculine, but \mr{पगारवाढ} follows \mr{वाढ} and is feminine.

\begin{cousinweb}
\textbf{\mr{पगारवाढ}} (\emph{pagārvāḍh}) is a Marathi compound of \mr{पगार}, \textquotedblleft{}salary\textquotedblright{}, a loan from Portuguese \emph{pagar}, \textquotedblleft{}to pay\textquotedblright{}, and \mr{वाढ}, \textquotedblleft{}growth\textquotedblright{}, from \mr{वाढणे}.
\end{cousinweb}

\begin{grammarlens}[title={What you've built}]
One more word for describing what you lived and felt.
\end{grammarlens}

\subsection*{Your turn}
\begin{itemize}
\raggedright
  \item \emph{Say it:} \emph{pagārvāḍh}
  \item \emph{Say it:} \emph{pagārvāḍh}, once more
  \item \emph{Say it:} \emph{yā varṣī sagḷyānnā cāṅglī pagārvāḍh miḷālī.} — Everyone got a good pay rise this year.
  \item \emph{Recall:} say the Marathi for daily wages, then the Marathi for an honorarium, then say \emph{pagārvāḍh} again
\end{itemize}

\begin{tcolorbox}[breakable,colback=teal!4,colframe=teal!35!black,title={Before you move on}]
What does \mr{पगारवाढ} mean? (\textquotedblleft{}A pay rise\textquotedblright{}.) Say it once more.
\end{tcolorbox}
\section[\texorpdfstring{\mr{कमावणे} (kamāvṇe)}{कमावणे (kamāvṇe)}]{\mr{कमावणे} (kamāvṇe) — to earn}

\label{lesson:MR-C508-kamavne}



\begin{quote}
Before the new one: say the Marathi for an honorarium, then the Marathi for a pay rise.
\end{quote}

\subsection*{You'll want to know: \mr{कमावणे}}
\textbf{\mr{कमावणे}} — \emph{kamāvṇe} — \textquotedblleft{}to earn\textquotedblright{}.

\begin{quote}
\mr{अरुणने} \mr{पहिल्या} \mr{नोकरीत} \mr{थोडे} \mr{पैसे} \mr{कमावले}.
\end{quote}

\begin{quote}
\emph{aruṇne pahilyā nokrīt thoḍe paise kamāvle.}
\end{quote}

\begin{quote}
Arun earned a little money in his first job.
\end{quote}

Past, with -\mr{ने} (he …): \mr{त्याने} \mr{पैसे} \mr{कमावले}, \emph{tyāne paise kamāvle}.

With \mr{पैसे} the past is plural: \mr{पैसे} \mr{कमावले}.

\begin{cousinweb}
\textbf{\mr{कमावणे}} (\emph{kamāvṇe}) is inherited: CDIAL lists Marathi \emph{kamāviṇẽ}, \textquotedblleft{}works, earns\textquotedblright{}, under \emph{*karmāpayati} (CDIAL 2897), built on Sanskrit \emph{karman}, \textquotedblleft{}work\textquotedblright{}.
\end{cousinweb}

\begin{grammarlens}[title={What you've built}]
One more word for describing what you lived and felt.
\end{grammarlens}

\subsection*{Your turn}
\begin{itemize}
\raggedright
  \item \emph{Say it:} \emph{kamāvṇe}
  \item \emph{Say it:} \emph{kamāvṇe}, once more
  \item \emph{Say it:} \emph{aruṇne pahilyā nokrīt thoḍe paise kamāvle.} — Arun earned a little money in his first job.
  \item \emph{Recall:} say the Marathi for an honorarium, then the Marathi for a pay rise, then say \emph{kamāvṇe} again
\end{itemize}

\begin{tcolorbox}[breakable,colback=teal!4,colframe=teal!35!black,title={Before you move on}]
What does \mr{कमावणे} mean? (\textquotedblleft{}To earn\textquotedblright{}.) Say it once more.
\end{tcolorbox}
